\section{Nanomaterial-Enabled Marine Antifouling Coatings}
\subsection{Metal and metal-oxide systems}
ZnO, TiO$_2$, Ag, CuO, Fe$_2$O$_3$, and mixed oxides are represented across the corpus. Reported mechanisms include antimicrobial activity, photocatalysis, altered wettability, and surface texture. For a transition-metal-doped oxide, a mechanistic interpretation should distinguish localized impurity or defect states from a simple increase in surface hydroxyl groups: dopant-derived $3d$ states can introduce sub-band-gap levels, alter carrier transfer, and change the probability that photogenerated electrons and holes reach the interface. Excess dopant, however, can create recombination centers and reduce photocatalytic efficiency under solar or marine illumination. The present corpus does not contain the cited DFT records needed to assign a specific dopant state, band position, or optimal concentration, so those quantities remain NR and require direct electronic-structure verification. Biosynthesized ZnO caused membrane damage and protein leakage in marine biofilm-forming bacteria in one study \cite{dharmaraj_protein_2021}; chitosan--ZnO and TiO$_2$-containing systems similarly combine matrix and particle functions \cite{al-naamani_chitosan-zinc_2017,alarif_green_2023}. These results support activity under reported test conditions, not a general claim of low toxicity.

\subsection{Carbon nanomaterials and hybrids}
CNTs, graphene, graphene oxide, carbon dots, and multi-component hybrids are used to modify strength, conductivity, photocatalysis, texture, and antimicrobial response. Dispersion is governed by matrix--filler interfacial energy, thermodynamic miscibility, polymer-chain adsorption, and the competition between filler--filler and filler--matrix interactions. In non-polar PDMS, poor compatibility can promote agglomeration even when nominal loading is low; agglomeration changes stress transfer, exposed active area, and release pathways \cite{cavas_reinforcement_2018}. The corpus does not report a common atomistic dispersion metric across studies, so loading values cannot be ranked independently of processing and matrix chemistry. Hybrid CuO--TiO$_2$--carbon-dot systems and carbon-dot--graphene systems extend the active-surface concept to aquaculture nets, where flexibility, abrasion, light availability, and release differ from hull coatings \cite{ashraf_photocatalytic_2024}.

\subsection{Evidence boundary and environmental trade-off}
Nanomaterial studies commonly demonstrate material characterization and short immersion or laboratory biological endpoints, whereas long-term wear, leaching, ecotoxicity, and full-scale hydrodynamics are less consistently reported. Metal release and impacts on sessile organisms are explicit concerns for copper and zinc systems \cite{miller_nano_2020}. The defensible design rule is to treat nanoparticle loading, dispersion, matrix retention, service abrasion, and end-of-life release as a coupled performance variable.

\subsection{Mechanistic depth: dopant state, band structure, ROS, and dispersion}
For photocatalytic systems, the label ``carbon-doped'' or ``hybrid oxide'' is not a mechanism by itself. A defensible interpretation should distinguish lattice substitution, interstitial incorporation, surface adsorption, and a separate heterojunction phase. The marine-extract TiO$_2$ study used XRD, SEM, EDS, XPS, and optical measurements, providing a useful starting characterization set \cite{alarif_green_2023}; however, composition alone does not establish the dopant state or the direction of charge transfer. Future reports should pair high-resolution XPS peak fitting with an appropriate reference, Raman or complementary structural mapping, and, where feasible, X-ray absorption or electron-paramagnetic-resonance evidence. The reported oxidation state, defect population, and dopant location should be stated separately from the nominal precursor composition.

Band-structure claims should likewise be tied to measurements. UV--visible diffuse reflectance can estimate an optical absorption edge, but a Tauc construction is sensitive to the assumed transition and should not be presented as a complete band diagram. Valence-band XPS, work-function or Kelvin-probe measurements, and Mott--Schottky analysis can constrain band positions and carrier type; photoluminescence quenching or time-resolved photoluminescence can indicate changed recombination, but neither alone proves higher photocatalytic quantum efficiency. For CuO--TiO$_2$--carbon-dot systems, action spectra and dark/light controls should be reported alongside wavelength-resolved irradiance so that apparent antifouling activity is not attributed to photocatalysis when it could arise from heating, dissolved copper, or an organic biocide \cite{ashraf_photocatalytic_2024}.

ROS pathways should be resolved with parallel scavenger and product controls. Hydroxyl-radical, superoxide, and singlet-oxygen probes can identify plausible species, but probe selectivity, seawater chloride, carbonate alkalinity, dissolved organic matter, and light attenuation can alter the result. A stronger protocol combines radical detection with oxygen demand, peroxide or oxidation-product measurements, dark controls, catalyst-free controls, and organism-relevant exposure tests. Membrane damage or protein leakage is evidence of a biological outcome, as reported for biosynthesized ZnO, but it does not by itself identify the responsible ROS or exclude dissolved-ion toxicity \cite{dharmaraj_protein_2021}.

Dispersion should be quantified rather than described as ``uniform.'' Recommended descriptors include particle-size distributions and agglomerate area fraction from statistically sampled cross-sectional SEM or TEM images, spatial Raman or EDS maps, coating-thickness-normalized loading, and an explicitly reported coefficient of variation across fields of view. For liquid precursors, report hydrodynamic size, polydispersity, and zeta potential before incorporation; for cured coatings, repeat the spatial analysis after immersion and abrasion. This is important because increased filler loading can improve reinforcement until agglomeration changes stress transfer and surface exposure, a boundary condition observed for carbon-filled PDMS \cite{cavas_reinforcement_2018}.

Supplementary defect-physics literature emphasizes that sub-band-gap traps and polaronic charge localization can change carrier lifetimes and interfacial reaction pathways, but these mechanisms require direct spectroscopy rather than inference from a broadened absorption edge alone \cite{arxiv220102808,arxiv230506708}. Recommended validation templates are valence-band XPS, ultraviolet photoelectron spectroscopy, Mott--Schottky impedance, photoluminescence lifetime, and action-spectrum measurements; these should be interpreted jointly rather than as interchangeable proxies.

\subsection{Standardized release-testing framework}
Environmental trade-offs should be evaluated using a tiered protocol rather than a single static leachate measurement. Tier 1 is a closed-bottle seawater immersion with known exposed area, coating mass, salinity, pH, temperature, dissolved oxygen, light condition, and water-to-area ratio. Samples should be collected over a predeclared time series and separated into dissolved, colloidal, and particulate fractions using validated filtration or ultrafiltration. Dissolved metals and major elements can be measured by ICP-MS or ICP-OES; particle number and size should be measured with single-particle ICP-MS, electron microscopy, or an orthogonal particle-sizing method where applicable. Report both area-normalized release and a mass balance against the remaining coating.

Tier 2 adds realistic stressors: reciprocating abrasion or calibrated waterjet exposure, wet--dry cycling, UV exposure, and representative flow. The protocol should record coating mass loss, roughness, gloss or wettability change, adhesion, and antifouling activity before and after stressing. Tier 3 uses a mesocosm or field panel with matched unfilled-matrix, neat-substrate, and positive-control panels, paired water and sediment or biofilm sampling, and recovery of detached solids. Each tier should retain field blanks, procedural blanks, matrix spikes, duplicate samples, and particle-loss checks. Ecotoxicity should use at least one primary producer or algae, one invertebrate or larval endpoint, and a microbial endpoint selected for the receiving environment, with dissolved-ion and particle-only controls where they can be separated. The output should be a concentration--time profile and an uncertainty statement, not simply ``no visible effect.''

\section{Bioinspired, Biobased, and Non-Toxic Surface Strategies}
\subsection{Hierarchical, slippery, and hydrated interfaces}
The corpus includes lotus-like, mucus-like, ciliary, superhydrophobic, superoleophobic, lubricant-infused, zwitterionic, heparin, hydrogel, and amphiphilic surfaces. Their common idea is to reduce the effective interaction between an organism and the substrate through low adhesion, hydration, mobility, or a structured interface. Gradient zwitterionic coatings produced by initiated chemical-vapour-deposition are a promising candidate architecture for bacterial anti-adhesion, but the required external study record is not present in the available evidence base; organism-specific performance, mechanical fragility, and long-term shear stability are therefore NR here. Lubricant-infused surfaces can be effective in laboratory or immersion tests, but lubricant retention, replenishment, contamination, and mechanical damage determine practical value \cite{basu_green_2020,liu_antifouling_2023}.

\subsection{Biobased chemistry}
Chitosan, capsaicin, marine natural products, diketopiperazines, xanthones, deep eutectic solvents, and other bio-derived compounds broaden the design space \cite{grant_towards_2022}. Biobased origin does not establish biodegradability in service or non-toxicity to non-target organisms; the complete formulation, exposure route, dose, and release pathway remain necessary.

\subsection{Realistic-condition limitation}
Bioinspired interfaces often show strong property-level or coupon-level results but can lose function under abrasion, pressure, biofilm contamination, UV exposure, or cleaning. Their most credible near-term role is as part of a maintenance-compatible system, with durability and adhesion measured alongside settlement.

\subsection{Lubricant-infused design guidelines and durability evidence}
The available studies support design rules, but not a universal optimum. The substrate should provide a connected retention structure and chemical affinity for the lubricant; the relevant roughness is therefore a texture-dependent design variable, not a single target $R_a$. One cerium-stearate study reported structured-surface average roughness of approximately $1.3\pm0.2~\mu$m for one morphology and found that a nanosheet-containing architecture retained lubricant more effectively during 40 days of artificial-seawater immersion than a less-retentive morphology \cite{yang_facile_2021}. A separate epoxy--TiO$_2$ SLIPS study examined silicone-oil viscosities from approximately 10 to 500 mPa s and used a coating roughness of about $7.24~\mu$m; the lower-viscosity configuration lost liquid-repellency behavior under scouring, while the 500 mPa s configuration showed the stronger retention response in the reported abrasion comparison \cite{yang_construction_2023}. These are study-specific operating windows, not universal viscosity or roughness prescriptions.

The practical guideline is to select viscosity jointly with texture, capillary retention, expected shear, and repairability: low viscosity can improve spreading but increase drainage and scouring loss, whereas high viscosity can improve retention but increase hysteresis and reduce self-healing or mobility. Report lubricant viscosity at test temperature, lubricant volume or mass uptake, texture dimensions and roughness distribution, contact-angle hysteresis, and lubricant mass or chemical tracer before and after immersion, flow, abrasion, and cleaning. Basu et al. found that UV pretreatment improved methyl-oleate infusion, with swelling ratios of $1.3\pm0.1$ for untreated and $1.9\pm0.4$ for UV-treated methyl-oleate-infused PDMS, and evaluated mussel response for two weeks; this supports an infusion/retention strategy but is not long-term ship-service evidence \cite{basu_green_2020}. Liu et al. reported stable sliding-angle behavior after 30 days for a lubricant-free slippery architecture, which should be distinguished from lubricant-retention evidence because no free lubricant reservoir is being maintained \cite{liu_antifouling_2023}.

Capillary-bridge formation provides an additional failure mode: an excess lubricant meniscus can form a bridge between a settling organism and the textured substrate, increasing adhesion rather than reducing it under some low-shear configurations \cite{arxiv201208619}. This mechanism should be checked alongside sliding-angle and adhesion measurements.

Across these studies, long-term immersion and abrasion data remain sparse and non-comparable. The minimum durability package for a marine SLIPS should therefore include repeated immersion with gravimetric or chemical lubricant-loss measurement, controlled flow or waterjet exposure, standardized abrasion cycles, post-stress sliding-angle and adhesion tests, coating integrity, and a matched release/ecotoxicity assessment. A surface that retains low contact-angle hysteresis but loses lubricant or changes chemistry should not be classified as durable without the mass balance.
